\section{Conclusion and Future Work}\label{sec:conclusion}

We applied Gruber's translation approach to the memory of LLM assistants: a shared ontology,
published and versioned like a software package, with thin per-engine adapters for write-in and
read-out. On a gold-first benchmark of multi-session personas and six cross-engine pairs of open
models, the ontology store carries memory between engines at a fraction of the tokens of a
transcript, and the remaining loss sits in extraction by the source engine. On long histories
the gold store reaches \PFM \res{pfm-oracle-large} where the transcript reaches
\res{pfm-b-large}, at \res{tok-ratio-oracle-b-large}$\times$ fewer tokens; stores written by
7--14B engines reach \res{pfm-a2-large}, but a store written once by a stronger model reaches
\res{pfm-a2x-opus-large} with the same small readers, matching a frontier model that reads the
whole transcript with \res{ptok-ratio-a2x-b-large}$\times$ fewer prompt tokens.

Future work: extraction by the conversational engine itself as models improve; hosted engines'
native memory exports as a fourth baseline; paraphrased and real dialogue; the regeneration
class of the taxonomy (moving retrieval corpora and re-embedding them); and consent scopes enforced
at read-out across organisations. The ontology, registry artifacts, benchmark, prompts and all run
outputs are available at \repo.
